\section{Quasi-determinization of separated finite-state automata}

We now describe an inductive construction that, given a $k$-separated finite-state
automaton $\mathcal{A} = \langle \Sigma, Q, X, S, F, \Delta \rangle$, produces a
$k$-separated quasi-deterministic finite-state automaton $\mathcal{A}'$ such that
$\mathcal{L}_{\mathcal{A}'} = \mathcal{L}_\mathcal{A}$.
The start state $s$ of the $k$-separated quasi-deterministic finite-state automaton
$\mathcal{A}'$ is defined as
\begin{equation*}
  \langle s^1, s^2, \ldots, s^k \rangle
\end{equation*}
where $s^i = \hat{\Delta}_\epsilon^i(S)$ for $i \in \{ 1, 2, \ldots, k \}$.
The base $k$-separated quasi-deterministic finite-state automaton
$\mathcal{A}'^{(1)}$ is defined as
\begin{equation*}
  \langle \Sigma, Q'^{(1)}, 2^X, s, F'^{(1)}, \tau^{(1)}, \rho^{(1)}, \eta^{(1)},
  \emptyset \rangle
\end{equation*}
where $Q_i'^{(1)} = \{ s^i \}$ and $F_i'^{(1)} = \tau_i^{(1)} = \rho_i^{(1)} =
\eta_i^{(1)} = \emptyset$ for $i \in \{ 1, 2, \ldots, k \}$.
The base set of states $R^{(1)}$ is defined as $\{ s \}$.
Let us assume that $\mathcal{A}'^{(n)}$ and $R^{(n)}$ are constructed for $n > 0$.
The set of states introduced at step $n$ is $R_n = R^{(n)} \setminus R^{(n-1)}$ where
$R^{(0)} = \emptyset$.
The next set of final quasi-states $F_i'^{(n+1)}$ is defined as
\begin{align*}
  F_i'^{(n)} \cup \{ q^i \mid \exists q \in R_n(q^i \cap F \neq \emptyset) \}
\end{align*}
for $i \in \{ 1, 2, \ldots, k \}$.
The next quasi-state transition function $\tau_i^{(n+1)}$ is defined as
\begin{align*}
  \tau_i^{(n)} \cup \{ \langle q^i, \sigma,
  \hat{\Delta}_\epsilon^i(\hat{\Delta}_f(q^i, \{ \sigma \})) \rangle \mid \exists q
  \in R_n \; \exists \sigma \in \Sigma \}
\end{align*}
for $i \in \{ 1, 2, \ldots, k \}$.
The next quasi-state reduction function $\rho_i^{(n+1)}$ is defined as
\begin{align*}
  \rho_i^{(n)} \cup \{ \langle q, q \cap X \rangle \mid \exists q \in
  rng(\tau_i^{(n+1)}) \setminus rng(\tau_i^{(n)}) \}
\end{align*}
for $i \in \{ 1, 2, \ldots, k \}$.
The next quasi-state extension function $\eta_i^{(n+1)}$ is defined as
\begin{align*}
  \eta_i^{(n)} \cup
    \{ &\langle t_i(q, \sigma), x(q, \sigma), t_i(q, \sigma) \cup
    \hat{\Delta}_\epsilon^i(x(q, \sigma)) \rangle \mid \exists q \in R_n \; \exists
    \sigma \in \Sigma \}
\end{align*}
where $t_i(q, \sigma) = \tau_i^{(n+1)}(q^i, \sigma)$ and $x(q, \sigma) =
\delta_x^{(n+1)}(q, \sigma)$ for $i \in \{ 1, 2, \ldots, k \}$.
The next separator extension function $\eta_x^{(n+1)}$ is defined as
\begin{align*}
  \eta_x^{(n)} \cup \bigcup_{i=1}^{k-1}
    \{ \langle \bigcup_{j=1}^i r_j(q, \sigma), r_{i+1}(q, \sigma),
    \bigcup_{j=1}^{i+1} r_j(q, \sigma)
    \rangle \mid \exists q \in R_n \; \exists \sigma \in \Sigma \}
\end{align*}
where $r_i(q, \sigma) = \rho_i^{(n+1)}(\tau_i^{(n+1)}(q^i, \sigma))$ for $i \in \{
1, 2, \ldots, k \}$.
The next set of states $Q_i'^{(n+1)}$ is defined as
\begin{equation*}
  Q_i'^{(n)} \cup rng(\tau_i^{(n+1)}) \cup rng(\eta_i^{(n+1)})
\end{equation*}
for $i \in \{ 1, 2, \ldots, k \}$.
The next $k$-separated quasi-deterministic finite-state automaton
$\mathcal{A}'^{(n+1)}$ is defined as
\begin{equation*}
  \langle \Sigma, Q'^{(n+1)}, 2^X, s, F'^{(n+1)}, \tau^{(n+1)}, \rho^{(n+1)},
  \eta^{(n+1)}, {\eta_x}^{(n+1)} \rangle.
\end{equation*}
The next set of states $R^{(n+1)}$ is defined as
\begin{equation*}
  R^{(n)} \cup \{ \delta^{(n+1)}(q, \sigma) \mid \exists q \in R_n \; \exists
  \sigma \in \Sigma \}.
\end{equation*}

\begin{proposition}
  Let $\mathcal{A}'^{(n)}$ be the $k$-separated quasi-deterministic finite-state
  automaton derived from the $k$-separated finite-state automaton $\mathcal{A} =
  \langle \Sigma, Q, X, S, F, \Delta \rangle$ by applying the construction described
  above for $n > 0$. Then $\tau_i^{(n)}, \rho_i^{(n)}, \eta_i^{(n)}$ for $i \in \{ 1,
  2, \ldots, k \}$ and $\eta_x^{(n)}$ are functions.
\end{proposition}

\begin{proof}
  Since $\hat{\Delta}_f$ and $\hat{\Delta}_\epsilon^i$ are functions, then
  $\tau_i^{(n)}$ is also a function for $i \in \{ 1, 2, \ldots, k \}$. By
  construction, $\rho_i^{(n)}$ is a function for $i \in \{ 1, 2, \ldots, k \}$. Since
  $\tau_i^{(n)}$ and $\rho_i^{(n)}$ are functions for $i \in \{ 1, 2, \ldots, k \}$,
  then $\eta_x^{(n)}$ is also a function. Since $\tau_i^{(n)}$ is a function for $i
  \in \{ 1, 2, \ldots, k \}$ and $\eta_x^{(n)}$ is a function, then $\eta_i^{(n)}$
  is also a function for $i \in \{ 1, 2, \ldots, k \}$.
\end{proof}

\begin{proposition}
  Let $\mathcal{A}'^{(n)}$ be the $k$-separated quasi-deterministic finite-state
  automaton derived from the $k$-separated finite-state automaton $\mathcal{A} =
  \langle \Sigma, Q, X, S, F, \Delta \rangle$ by applying the construction described
  above for $n > 0$. Then $dom(\rho_i^{(n)}) = rng(\tau_i^{(n)})$ for $i \in \{ 1, 2,
  \ldots, k \}$.
\end{proposition}

\begin{proof}
  We show that $dom(\rho_i^{(n)}) = rng(\tau_i^{(n)})$ for $i \in \{ 1, 2, \ldots, k
  \}$ by induction on $n$. Since $\rho_i^{(1)} = \tau_i^{(1)} = \emptyset$ for $i \in
  \{ 1, 2, \ldots, k \}$, the statement holds for $n = 1$. Let us assume that
  $dom(\rho_i^{(n)}) = rng(\tau_i^{(n)})$ for $n > 0$ and $i \in \{ 1, 2, \ldots, k
  \}$. Then
  \begin{align*}
    dom(\rho_i^{(n+1)}) & = dom(\rho_i^{(n)}) \cup (rng(\tau_i^{(n+1)}) \setminus
    rng(\tau_i^{(n)}))\\
                        & = rng(\tau_i^{(n)}) \cup (rng(\tau_i^{(n+1)}) \setminus
    rng(\tau_i^{(n)}))\\
                        & = rng(\tau_i^{(n+1)}).
  \end{align*}
\end{proof}

\begin{lemma}\label{lem:quasidet_completion}
  Let $\mathcal{A}'^{(n)}$ and $R^{(n)}$ be respectively the $k$-separated
  quasi-deterministic finite-state automaton and the set of states derived from the
  $k$-separated finite-state automaton $\mathcal{A} = \langle \Sigma, Q, X, S, F,
  \Delta \rangle$ by applying the construction described above for $n > 0$. There
  exists $m > 0$ such that $R^{(n)} = R^{(n+1)}$ and $\mathcal{A}'^{(n)} =
  \mathcal{A}'^{(n+1)}$ for $n > m$.
\end{lemma}

\begin{proof}
  By construction, $R^{(n)} \subseteq R^{(n+1)} \subseteq
  \prod_{i=1}^k{2^{\mathcal{C}(Q_i)}}$ for $n > 0$. Therefore, there exists $m \leq
  |\prod_{i=1}^k{2^{\mathcal{C}(Q_i)}}|$ such that $R^{(m)} = R^{(m+1)}$ and thus
  $R_{m+1} = \emptyset$. In that case, $\mathcal{A}'^{(m+1)} = \mathcal{A}'^{(m+2)}$
  and $R^{(m+1)} = R^{(m+2)}$. It follows that $R^{(n)} = R^{(n+1)}$ and
  $\mathcal{A}'^{(n)} = \mathcal{A}'^{(n+1)}$ for $n > m$.
\end{proof}

For the smallest $n > 0$ such that $\mathcal{A}'^{(n)} = \mathcal{A}'^{(n+1)}$ we
define $\mathcal{A}' = \mathcal{A}'^{(n)}$ and $R = R^{(n)}$.

By construction, $s \in R$, and whenever a state $q$ is introduced in $R_n$, the
entries for $q$ and every $\sigma \in \Sigma$ are added at step $n + 1$ and $\delta(q,
\sigma) \in R^{(n+1)}$. Therefore, $\delta^*(s, w)$ is defined for every $w \in
\Sigma^*$, and $R = \{ \delta^*(s, w) \mid w \in \Sigma^* \}$.
For $w \in \Sigma^*$ let
\begin{equation*}
  Reach(w) = \{ q \in Q \mid \exists s' \in S (\langle s', w, q \rangle \in \Delta^*)
  \}
\end{equation*}
be the set of states reachable from an initial state of $\mathcal{A}$ by reading $w$.
The following lemma shows that the $i$-th quasi-state of $\delta^*(s, w)$ is exactly
the part of $Reach(w)$ that lies in the component $\mathcal{C}(Q_i)$.

\begin{lemma}\label{lem:quasidet_invariant}
  Let $\mathcal{A}' = \langle \Sigma, Q', 2^X, s, F', \tau, \rho, \eta, \eta_x
  \rangle$ be the $k$-separated quasi-deterministic finite-state automaton derived
  from the $k$-separated finite-state automaton $\mathcal{A} = \langle \Sigma, Q, X,
  S, F, \Delta \rangle$ by applying the construction described above. Then
  \begin{equation*}
    \delta^*(s, w)^i = Reach(w) \cap \mathcal{C}(Q_i)
  \end{equation*}
  for $w \in \Sigma^*$ and $i \in \{ 1, 2, \ldots, k \}$.
\end{lemma}

\begin{proof}
  We prove the statement by induction on the length of $w$. If $|w| = 0$, then $w =
  \epsilon$ and
  \begin{align*}
    \delta^*(s, \epsilon)^i = s^i = \hat{\Delta}_\epsilon^i(S) &= \{ q \in
    \mathcal{C}(Q_i) \mid \exists s' \in S (\langle s', \epsilon, q \rangle \in
    \Delta^*) \}\\
    &= Reach(\epsilon) \cap \mathcal{C}(Q_i).
  \end{align*}
  Otherwise, $w = u\sigma$ for $u \in \Sigma^*$ and $\sigma \in \Sigma$. Let $p =
  \delta^*(s, u)$. By the inductive hypothesis, $p^r = Reach(u) \cap
  \mathcal{C}(Q_r)$ for $r \in \{ 1, 2, \ldots, k \}$. By construction,
  \begin{align*}
    \delta^*(s, w)^i = \delta_i(p, \sigma) &= \tau_i(p^i, \sigma) \cup
    \hat{\Delta}_\epsilon^i(\delta_x(p, \sigma)),\\
    \tau_r(p^r, \sigma) &= \hat{\Delta}_\epsilon^r(\hat{\Delta}_f(p^r, \{ \sigma
    \})),\\
    \delta_x(p, \sigma) &= \bigcup_{r=1}^k (\tau_r(p^r, \sigma) \cap X).
  \end{align*}
  Observe that $\tau_r(p^r, \sigma) \subseteq Reach(w)$ for $r \in \{ 1, 2, \ldots, k
  \}$. Indeed, if $q \in \tau_r(p^r, \sigma)$, then there exist $q' \in p^r \subseteq
  Reach(u)$ and $q'' \in Q$ such that $\langle q', \sigma, q'' \rangle \in \Delta$ and
  $\langle q'', \epsilon, q \rangle \in \Delta^*$, and thus $q \in Reach(u\sigma)$.

  First, let $q \in \delta^*(s, w)^i$. Then $q \in \mathcal{C}(Q_i)$. If $q \in
  \tau_i(p^i, \sigma)$, then $q \in Reach(w)$ by the observation above. Otherwise, $q
  \in \hat{\Delta}_\epsilon^i(\delta_x(p, \sigma))$ and there exist $r \in \{ 1, 2,
  \ldots, k \}$ and $x \in \tau_r(p^r, \sigma) \cap X$ such that $\langle x,
  \epsilon, q \rangle \in \Delta^*$. Since $x \in Reach(w)$, also $q \in Reach(w)$.

  Conversely, let $q \in Reach(w) \cap \mathcal{C}(Q_i)$. Consider a path with label
  $w$ from an initial state to $q$, and the last transition on it labelled with
  $\sigma$. There exist $s' \in S$ and $q', q'' \in Q$ such that $\langle s', u, q'
  \rangle \in \Delta^*$, $\langle q', \sigma, q'' \rangle \in \Delta$ and $\langle
  q'', \epsilon, q \rangle \in \Delta^*$. In particular, $q' \in Reach(u)$. If $q'
  \in \mathcal{C}(Q_i)$, then $q' \in p^i$ and thus $q \in \tau_i(p^i, \sigma)
  \subseteq \delta^*(s, w)^i$. Otherwise, $q' \in Q_r$ for some $r \neq i$. Since
  $\mathcal{A}$ is $k$-separated, $q'' \not\in Q_j$ for $j \neq r$, and thus $q'' \in
  \mathcal{C}(Q_r)$. Since $q'' \in \mathcal{C}(Q_r)$, $q \in \mathcal{C}(Q_i)$ and $r
  \neq i$, by \Cref{prop:separator_path} there exist $w', w'' \in \Sigma^*$ and $x \in
  X$ such that $\epsilon = w'w''$ and $\langle q'', w', x \rangle, \langle x, w'', q
  \rangle \in \Delta^*$. Thus, $w' = w'' = \epsilon$. Since $q' \in Reach(u) \cap
  \mathcal{C}(Q_r) = p^r$ and $x \in X \subseteq \mathcal{C}(Q_r)$, we have $x \in
  \tau_r(p^r, \sigma) \cap X \subseteq \delta_x(p, \sigma)$. Therefore, $q \in
  \hat{\Delta}_\epsilon^i(\delta_x(p, \sigma)) \subseteq \delta^*(s, w)^i$.
\end{proof}

In particular, $\delta^*(s, w)^i \cap X = \delta^*(s, w)^j \cap X = Reach(w) \cap X$
for $i, j \in \{ 1, 2, \ldots, k \}$: every reachable separator state belongs to all
quasi-states. The separator transition function $\delta_x$ keeps the components
synchronized in this way.

\begin{lemma}\label{lem:quasidet_inclusion_right}
  Let $\mathcal{A}' = \langle \Sigma, Q', 2^X, s, F', \tau, \rho, \eta, \eta_x
  \rangle$ be the $k$-separated quasi-deterministic finite-state automaton derived
  from the $k$-separated finite-state automaton $\mathcal{A} = \langle \Sigma, Q, X,
  S, F, \Delta \rangle$ by applying the construction described above. Then
  $\mathcal{L}_\mathcal{A} \subseteq \mathcal{L}_{\mathcal{A}'}$.
\end{lemma}

\begin{proof}
  Let $w \in \mathcal{L}_\mathcal{A}$. There exist $s' \in S$ and $f \in F$ such that
  $\langle s', w, f \rangle \in \Delta^*$. Thus, $f \in Reach(w)$. Since $Q = X \cup
  \bigcup_{i=1}^k Q_i$, there exists $i \in \{ 1, 2, \ldots, k \}$ such that $f \in
  \mathcal{C}(Q_i)$. By \Cref{lem:quasidet_invariant}, $f \in \delta^*(s, w)^i$, and
  thus $\delta^*(s, w)^i \cap F \neq \emptyset$. Since $\delta^*(s, w) \in R$, by
  construction $\delta^*(s, w)^i \in F_i'$. Therefore, $w \in
  \mathcal{L}_{\mathcal{A}'}$.
\end{proof}

\begin{lemma}\label{lem:quasidet_inclusion_left}
  Let $\mathcal{A}' = \langle \Sigma, Q', 2^X, s, F', \tau, \rho, \eta, \eta_x
  \rangle$ be the $k$-separated quasi-deterministic finite-state automaton derived
  from the $k$-separated finite-state automaton $\mathcal{A} = \langle \Sigma, Q, X,
  S, F, \Delta \rangle$ by applying the construction described above. Then
  $\mathcal{L}_{\mathcal{A}'} \subseteq \mathcal{L}_\mathcal{A}$.
\end{lemma}

\begin{proof}
  Let $w \in \mathcal{L}_{\mathcal{A}'}$. There exists $i \in \{ 1, 2, \ldots, k \}$
  such that $\delta^*(s, w)^i \in F_i'$. By construction, every element of $F_i'$
  intersects $F$. Thus, there exists $f \in \delta^*(s, w)^i \cap F$. By
  \Cref{lem:quasidet_invariant}, $f \in Reach(w)$, and thus there exists $s' \in S$
  such that $\langle s', w, f \rangle \in \Delta^*$. Therefore, $w \in
  \mathcal{L}_\mathcal{A}$.
\end{proof}

\begin{theorem}
  For every $k$-separated finite-state automaton $\mathcal{A}$ there exists a
  $k$-separated quasi-deterministic finite-state automaton $\mathcal{A}'$ such that
  $\mathcal{L}_{\mathcal{A}'} = \mathcal{L}_\mathcal{A}$.
\end{theorem}

\begin{proof}
  Follows from \Cref{lem:quasidet_completion,lem:quasidet_inclusion_right,lem:quasidet_inclusion_left}.
\end{proof}

The construction described above inspires \Cref{alg:quasi_det_sfsa} for
quasi-determinization of finite-state automata. We shall now prove its correctness
and provide a runtime complexity bound.
\pagebreak

  \SetAlgoNoLine
  \begin{algorithm}[H]
    \caption{Quasi-determinization of a k-separated finite-state automaton}
    \label{alg:quasi_det_sfsa}

    \DontPrintSemicolon
    \KwIn{$\langle \Sigma, Q, X, S, F, \Delta \rangle$}
    \KwOut{$\langle \Sigma, Q', s, F', \tau, \rho, \eta, \eta_x \rangle$}

    $S_\epsilon \gets \hat{\Delta}_\epsilon(S)$\;\nllabel{alg:qd:init_begin}

    \For {$i \in \{ 1, 2, \ldots, k \}$} {
      $P_i \gets Q_i \cup X$\;
      $s_i \gets S_\epsilon \cap P_i$\;
      $Q_i' \gets \{ s_i \}$\;
      $F_i' \gets \emptyset$, $\tau_i \gets \emptyset$, $\rho_i \gets \emptyset$,
      $\eta_i \gets \emptyset$\;
    }
    $\eta_x \gets \emptyset$\;

    $Next \gets \{ s \}$, $Visited \gets \emptyset$\;\nllabel{alg:qd:init_end}

    \While {$Next \neq \emptyset$} {
      $q \gets Extract(Next)$\;\nllabel{alg:qd:extract}
      $Visited \gets Visited \cup \{ q \}$\;\nllabel{alg:qd:visited}
      \For {$i \in \{ 1, 2, \ldots, k \}$} {
        \If {$q_i \cap F \neq \emptyset$} {
          $F_i' \gets F_i' \cup \{ q_i \}$\;\nllabel{alg:qd:final}
        }
      }

      \For {$\sigma \in \Sigma$} {
        \For {$i \in \{ 1, 2, \ldots, k \}$} {
          $t_i \gets \hat{\Delta}_\epsilon(\hat{\Delta}_f(q_i, \{
            \sigma \})) \cap P_i$\;\nllabel{alg:qd:t}
          $\tau_i \gets \tau_i \cup \{ \langle q_i, \sigma, t_i
            \rangle \}$\;
          $Q_i' \gets Q_i' \cup \{ t_i \}$\;
          $r_i \gets t_i \cap X$\;
          $\rho_i \gets \rho_i \cup \{ \langle t_i, r_i \rangle
            \}$\;\nllabel{alg:qd:rho}
        }

        $e_x \gets r_1$\;\nllabel{alg:qd:ex_begin}
        \For {$i \in \{ 2, 3, \ldots, k \}$} {
          $\eta_x \gets \eta_x \cup \{ \langle e_x, r_i, e_x \cup r_i \rangle \}$\;
          $e_x \gets e_x \cup r_i$\;\nllabel{alg:qd:ex_end}
        }

        \For {$i \in \{ 1, 2, \ldots, k \}$} {
          $e_i \gets t_i \cup (\hat{\Delta}_\epsilon(e_x) \cap P_i)$\;\nllabel{alg:qd:e}
          $\eta_i \gets \eta_i \cup \{ \langle t_i, e_x, e_i \rangle \}$\;
          $Q_i' \gets Q_i' \cup \{ e_i \}$\;\nllabel{alg:qd:e_end}
        }

        \If {$\langle e_1, e_2, \ldots, e_k \rangle \not\in Visited$} {
          $Next \gets Next \cup \{ \langle e_1, e_2, \ldots, e_k \rangle
            \}$\;\nllabel{alg:qd:next}
        }
      }
    }

    \Return $\langle \Sigma, Q', s, F', \tau, \rho, \eta, \eta_x \rangle$
  \end{algorithm}

  \begin{proposition}\label{prop:quasidet_algorithm}
    Let $\mathcal{A} = \langle \Sigma, Q, X, S, F, \Delta \rangle$ be a
    $k$-separated finite-state automaton. Let $\mathcal{A}_c$ be the
    quasi-deterministic $k$-separated finite-state automaton produced by
    \Cref{alg:quasi_det_sfsa} on input $\mathcal{A}$. Then every state of $R$ is
    extracted from $Next$ at line~\ref{alg:qd:extract} exactly once, no other state is
    extracted, and $\mathcal{A}_c = \mathcal{A}'$.
  \end{proposition}

  \begin{proof}
    A state is added to $Next$ only at line~\ref{alg:qd:next}, and only if it is not
    in $Visited$. Every extracted state is added to $Visited$ at
    line~\ref{alg:qd:visited} and is never removed from it. Therefore, no state is
    extracted twice.

    Consider the iteration of the while loop in which a state $q$ is extracted.
    Lines~\ref{alg:qd:t}--\ref{alg:qd:rho} compute, for each $\sigma \in \Sigma$,
    \begin{align*}
      t_i &= \hat{\Delta}_\epsilon(\hat{\Delta}_f(q_i, \{ \sigma \})) \cap P_i =
      \hat{\Delta}_\epsilon^i(\hat{\Delta}_f(q^i, \{ \sigma \})) = \tau_i(q^i,
      \sigma),\\
      r_i &= t_i \cap X = \rho_i(t_i).
    \end{align*}
    Lines~\ref{alg:qd:ex_begin}--\ref{alg:qd:ex_end} compute $e_x = \bigcup_{i=1}^k
    r_i = \delta_x(q, \sigma)$ and add the entries of $\eta_x$ for its prefixes, and
    lines~\ref{alg:qd:e}--\ref{alg:qd:e_end} compute $e_i = t_i \cup
    \hat{\Delta}_\epsilon^i(e_x) = \eta_i(t_i, e_x) = \delta_i(q, \sigma)$. Thus, the
    state considered at line~\ref{alg:qd:next} is $\delta(q, \sigma)$. Moreover, the
    entries added to $F'$, $\tau$, $\rho$, $\eta$, $\eta_x$ and $Q'$ in this iteration
    are exactly the entries that the construction adds for the state $q$.

    Every extracted state belongs to $R$, since $s \in R$ and $\delta(q, \sigma) \in R$
    for $q \in R$ and $\sigma \in \Sigma$. Conversely, every state $\delta^*(s, w) \in
    R$ is extracted. We show this by induction on the length of $w$. The state $s$ is
    extracted in the first iteration. If $\delta^*(s, u)$ is extracted, then in that
    iteration $\delta^*(s, u\sigma) = \delta(\delta^*(s, u), \sigma)$ is either already
    in $Visited$, in which case it has been extracted, or it is added to $Next$. Since
    the while loop terminates only when $Next = \emptyset$, it is extracted in a later
    iteration. The loop terminates, because each iteration extracts a different state of
    the finite set $R$.

    Therefore, \Cref{alg:quasi_det_sfsa} processes each state of $R$ exactly once and
    adds exactly the entries that the construction adds for the states of $R$. Since
    the construction adds entries only for states of $R$, $\mathcal{A}_c =
    \mathcal{A}'$.
  \end{proof}

  Let
  \begin{equation*}
    D = |\{ Reach(w) \mid w \in \Sigma^* \}|.
  \end{equation*}
  By definition, $D$ is the number of states reachable from the initial state of the
  deterministic finite-state automaton obtained from $\mathcal{A}$ by the construction
  of \Cref{thm:rabin_scott}.

  \begin{proposition}\label{prop:quasidet_iterations}
    Let $\mathcal{A} = \langle \Sigma, Q, X, S, F, \Delta \rangle$ be a
    $k$-separated finite-state automaton. Then $|R| = D \leq 2^{|Q|}$. Consequently,
    the while loop of \Cref{alg:quasi_det_sfsa} executes exactly $D$ iterations on
    input $\mathcal{A}$.
  \end{proposition}

  \begin{proof}
    By \Cref{lem:quasidet_invariant}, $\delta^*(s, w) = \langle Reach(w) \cap
    \mathcal{C}(Q_1), \ldots, Reach(w) \cap \mathcal{C}(Q_k) \rangle$ for $w \in
    \Sigma^*$. Since $Q = \bigcup_{i=1}^k \mathcal{C}(Q_i)$, also $Reach(w) =
    \bigcup_{i=1}^k \delta^*(s, w)^i$. Therefore, the map $\delta^*(s, w) \mapsto
    Reach(w)$ is a well-defined bijection between $R = \{ \delta^*(s, w) \mid w \in
    \Sigma^* \}$ and $\{ Reach(w) \mid w \in \Sigma^* \}$, and $|R| = D$. Clearly, $D
    \leq |2^Q| = 2^{|Q|}$. By \Cref{prop:quasidet_algorithm}, the while loop executes
    exactly $|R|$ iterations.
  \end{proof}

  \begin{proposition}\label{prop:quasidet_complexity}
    Let $\mathcal{A} = \langle \Sigma, Q, X, S, F, \Delta \rangle$ be a
    $k$-separated finite-state automaton. Then \Cref{alg:quasi_det_sfsa} completes on
    input $\mathcal{A}$ in $\mathcal{O}(k |\Sigma| |Q| D)$ time. In particular, it
    completes in $\mathcal{O}(k |\Sigma| |Q| 2^{|Q|})$ time.
  \end{proposition}

  \begin{proof}
    We assume that every operation on subsets of $Q$ (union, intersection, membership
    test, set-lifted transition and $\epsilon$-closure, and table look-up) takes
    $\mathcal{O}(|Q|)$ time. Look-ups and insertions of $k$-tuples of such subsets
    (lines~\ref{alg:qd:extract}, \ref{alg:qd:visited} and~\ref{alg:qd:next}) take
    $\mathcal{O}(k |Q|)$ time.
    Lines~\ref{alg:qd:init_begin}--\ref{alg:qd:init_end} are executed once and require
    $\mathcal{O}(k |Q|)$ time.
    By \Cref{prop:quasidet_iterations}, the body of the while loop is executed $D$
    times. In each execution, lines~\ref{alg:qd:extract}--\ref{alg:qd:final} require
    $\mathcal{O}(k |Q|)$ time. For each of the $|\Sigma|$ symbols,
    lines~\ref{alg:qd:t}--\ref{alg:qd:rho} are executed $k$ times,
    lines~\ref{alg:qd:ex_begin}--\ref{alg:qd:ex_end} at most $k$ times,
    lines~\ref{alg:qd:e}--\ref{alg:qd:e_end} $k$ times, and each of these executions
    requires $\mathcal{O}(|Q|)$ time; line~\ref{alg:qd:next} and the test before it
    require $\mathcal{O}(k |Q|)$ time. Thus, each iteration of the while loop requires
    $\mathcal{O}(k |\Sigma| |Q|)$ time. Therefore, the total runtime complexity bound
    for \Cref{alg:quasi_det_sfsa} is $\mathcal{O}(k |Q| + D k |\Sigma| |Q|) =
    \mathcal{O}(k |\Sigma| |Q| D)$. By \Cref{prop:quasidet_iterations}, $D \leq
    2^{|Q|}$.
  \end{proof}

  The size of the constructed device is bounded in terms of the components:
  $|Q_i'| \leq 2^{|\mathcal{C}(Q_i)|} = 2^{|Q_i| + |X|}$ for $i \in \{ 1, 2, \ldots, k
  \}$. The running time of \Cref{alg:quasi_det_sfsa}, however, depends on $D$, the
  number of states of the deterministic finite-state automaton that the construction
  avoids representing. This cannot be improved by a better analysis of the algorithm:
  the while loop ranges over the states of $\mathcal{A}'$, which are $k$-tuples of
  quasi-states, and by \Cref{prop:quasidet_iterations} there are exactly $D$ of them.
  The following example shows that $D$ can be exponentially larger than the sizes of
  the components.

  For $k > 1$, let $\mathcal{A}_k = \langle \Sigma_k, Q, X, S, F, \Delta \rangle$ where
  $\Sigma_k = \{ a_1, a_2, \ldots, a_k \}$, $Q = \{ x, q_1, q_2, \ldots, q_k \}$, $X =
  \{ x \}$, $Q_i = \{ q_i \}$ for $i \in \{ 1, 2, \ldots, k \}$, $S = \{ x \}$, $F = \{
  q_1, q_2, \ldots, q_k \}$, and
  \begin{equation*}
    \Delta = \{ \langle x, \epsilon, q_i \rangle \mid i \in \{ 1, 2, \ldots, k \} \}
    \cup \{ \langle q_i, a_j, q_i \rangle \mid i, j \in \{ 1, 2, \ldots, k \}, i \neq
    j \}.
  \end{equation*}
  The finite-state automaton $\mathcal{A}_k$ is $k$-separated with $|X| = |Q_i| = 1$,
  and it recognizes the words in which some symbol of $\Sigma_k$ does not occur. For $J
  \subseteq \{ 1, 2, \ldots, k \}$ let $w_J$ be a word in which exactly the symbols
  $a_j$ for $j \in J$ occur. Then $q_i \in Reach(w_J)$ if and only if $i \not\in J$.
  Thus, the sets $Reach(w_J)$ are pairwise different, and $D \geq 2^k$. On the other
  hand, $|Q_i'| \leq 2^{|Q_i| + |X|} = 4$, so $\mathcal{A}_k'$ is described by tables
  over at most $4k$ quasi-states. On input $\mathcal{A}_k$,
  \Cref{alg:quasi_det_sfsa} executes at least $2^k$ iterations. In particular, its
  running time cannot be bounded by a polynomial in $k |\Sigma_k| |Q| 2^{|X|}
  \sum_{i=1}^k 2^{|Q_i|} = 4k^3(k+1)$.

The tables of $\mathcal{A}'$ are indexed by component quasi-states and separator states
only. A global state is needed solely to determine which separator state $\delta_x(q,
\sigma)$ is passed to $\eta_i$. We now describe a variant of the construction that
passes every subset of $X$ instead, and therefore explores each component separately.

Let $\tau_i(P, \sigma) = \hat{\Delta}_\epsilon^i(\hat{\Delta}_f(P, \{ \sigma \}))$ and
$\eta_i(P, x) = P \cup \hat{\Delta}_\epsilon^i(x)$ for $P \subseteq \mathcal{C}(Q_i)$,
$\sigma \in \Sigma$ and $x \subseteq X$, as in the construction above. For $i \in \{ 1,
2, \ldots, k \}$ let $U_i$ be the smallest set such that
\begin{itemize}
  \item $s^i \in U_i$,
  \item $\eta_i(\tau_i(P, \sigma), x) \in U_i$ for $P \in U_i$, $\sigma \in \Sigma$ and
    $x \subseteq X$,
\end{itemize}
and let $T_i = \{ \tau_i(P, \sigma) \mid P \in U_i, \sigma \in \Sigma \}$. Since $U_i
\cup T_i \subseteq 2^{\mathcal{C}(Q_i)}$, we have $|U_i \cup T_i| \leq 2^{|Q_i| +
|X|}$. Let
\begin{equation*}
  \mathcal{A}'' = \langle \Sigma, Q'', 2^X, s, F'', \tau'', \rho'', \eta'', \eta_x''
  \rangle
\end{equation*}
where, for $i \in \{ 1, 2, \ldots, k \}$,
\begin{itemize}
  \item $Q_i'' = U_i \cup T_i$,
  \item $F_i'' = \{ P \in U_i \mid P \cap F \neq \emptyset \}$,
  \item $\tau_i''$ is the restriction of $\tau_i$ to $U_i \times \Sigma$,
  \item $\rho_i''(t) = t \cap X$ for $t \in T_i = rng(\tau_i'')$,
  \item $\eta_i''$ is the restriction of $\eta_i$ to $T_i \times 2^X$,
  \item $\eta_x''(x, y) = x \cup y$ for $x, y \subseteq X$.
\end{itemize}
By definition, $\tau_i''$ maps into $T_i$ and $\eta_i''$ maps into $U_i$, so
$\mathcal{A}''$ is a $k$-separated quasi-deterministic finite-state automaton.

\begin{lemma}\label{lem:quasidet_components}
  Let $\mathcal{A} = \langle \Sigma, Q, X, S, F, \Delta \rangle$ be a $k$-separated
  finite-state automaton. Then $\delta''^*(s, w)^i = Reach(w) \cap \mathcal{C}(Q_i) \in
  U_i$ for $w \in \Sigma^*$ and $i \in \{ 1, 2, \ldots, k \}$, and
  $\mathcal{L}_{\mathcal{A}''} = \mathcal{L}_\mathcal{A}$.
\end{lemma}

\begin{proof}
  Let $q \in \prod_{i=1}^k U_i$ and $\sigma \in \Sigma$. Then $\tau_i''(q^i, \sigma) =
  \tau_i(q^i, \sigma) \in T_i$, $\delta_x''(q, \sigma) = \bigcup_{i=1}^k (\tau_i(q^i,
  \sigma) \cap X) \subseteq X$, and $\delta_i''(q, \sigma) = \eta_i(\tau_i(q^i, \sigma),
  \delta_x''(q, \sigma)) \in U_i$. Thus, all values are defined and $\delta''(q, \sigma)
  \in \prod_{i=1}^k U_i$. They are given by the formulas used in the proof of
  \Cref{lem:quasidet_invariant}, and $s \in \prod_{i=1}^k U_i$. Therefore, the proof of
  \Cref{lem:quasidet_invariant} applies verbatim to $\mathcal{A}''$, and $\delta''^*(s,
  w)^i = Reach(w) \cap \mathcal{C}(Q_i) \in U_i$. Since $F_i''$ contains exactly the
  quasi-states in $U_i$ that intersect $F$, the proofs of
  \Cref{lem:quasidet_inclusion_right,lem:quasidet_inclusion_left} also apply verbatim,
  and $\mathcal{L}_{\mathcal{A}''} = \mathcal{L}_\mathcal{A}$.
\end{proof}

In particular, $\delta''^*(s, w) = \delta^*(s, w)$ for $w \in \Sigma^*$, so
$\mathcal{A}''$ behaves exactly like $\mathcal{A}'$. However, $U_i \cup T_i$ may contain
quasi-states that are not components of any reachable state, so $\mathcal{A}''$ may have
more quasi-states than $\mathcal{A}'$. Both are bounded by $2^{|Q_i| + |X|}$ in each
component. \Cref{alg:quasi_det_sfsa_components} computes $\mathcal{A}''$.

  \SetAlgoNoLine
  \begin{algorithm}[!ht]
    \caption{Per-component quasi-determinization of a k-separated finite-state
      automaton}
    \label{alg:quasi_det_sfsa_components}

    \DontPrintSemicolon
    \KwIn{$\langle \Sigma, Q, X, S, F, \Delta \rangle$}
    \KwOut{$\langle \Sigma, Q'', 2^X, s, F'', \tau'', \rho'', \eta'', \eta_x'' \rangle$}

    $S_\epsilon \gets \hat{\Delta}_\epsilon(S)$\;\nllabel{alg:qdc:init}
    $\eta_x'' \gets \{ \langle x, y, x \cup y \rangle \mid x, y \subseteq X
      \}$\;\nllabel{alg:qdc:etax}

    \For {$i \in \{ 1, 2, \ldots, k \}$} {
      $P_i \gets Q_i \cup X$\;\nllabel{alg:qdc:comp_begin}
      $s_i \gets S_\epsilon \cap P_i$\;
      $Q_i'' \gets \{ s_i \}$, $Next \gets \{ s_i \}$, $Visited \gets \emptyset$,
      $Extended \gets \emptyset$\;
      $F_i'' \gets \emptyset$, $\tau_i'' \gets \emptyset$, $\rho_i'' \gets \emptyset$,
      $\eta_i'' \gets \emptyset$\;\nllabel{alg:qdc:comp_init_end}

      \While {$Next \neq \emptyset$} {
        $P \gets Extract(Next)$\;\nllabel{alg:qdc:extract}
        $Visited \gets Visited \cup \{ P \}$\;
        \If {$P \cap F \neq \emptyset$} {
          $F_i'' \gets F_i'' \cup \{ P \}$\;\nllabel{alg:qdc:final}
        }
        \For {$\sigma \in \Sigma$} {
          $t \gets \hat{\Delta}_\epsilon(\hat{\Delta}_f(P, \{ \sigma \})) \cap
            P_i$\;\nllabel{alg:qdc:t}
          $\tau_i'' \gets \tau_i'' \cup \{ \langle P, \sigma, t \rangle \}$\;
          $\rho_i'' \gets \rho_i'' \cup \{ \langle t, t \cap X \rangle \}$\;
          $Q_i'' \gets Q_i'' \cup \{ t \}$\;\nllabel{alg:qdc:t_end}
          \If {$t \not\in Extended$} {\nllabel{alg:qdc:extended}
            $Extended \gets Extended \cup \{ t \}$\;
            \For {$x \subseteq X$} {
              $e \gets t \cup (\hat{\Delta}_\epsilon(x) \cap P_i)$\;\nllabel{alg:qdc:e}
              $\eta_i'' \gets \eta_i'' \cup \{ \langle t, x, e \rangle \}$\;
              $Q_i'' \gets Q_i'' \cup \{ e \}$\;
              \If {$e \not\in Visited$} {
                $Next \gets Next \cup \{ e \}$\;\nllabel{alg:qdc:e_next}
              }
            }
          }
        }
      }
    }

    \Return $\langle \Sigma, Q'', 2^X, s, F'', \tau'', \rho'', \eta'', \eta_x'' \rangle$
  \end{algorithm}

  \begin{proposition}\label{prop:quasidet_components_algorithm}
    Let $\mathcal{A} = \langle \Sigma, Q, X, S, F, \Delta \rangle$ be a
    $k$-separated finite-state automaton. Then \Cref{alg:quasi_det_sfsa_components}
    returns $\mathcal{A}''$ on input $\mathcal{A}$. Moreover, for each $i \in \{ 1, 2,
    \ldots, k \}$, every element of $U_i$ is extracted at line~\ref{alg:qdc:extract}
    exactly once, no other quasi-state is extracted, and every element of $T_i$ is
    added to $Extended$ at line~\ref{alg:qdc:extended} exactly once.
  \end{proposition}

  \begin{proof}
    Let $i \in \{ 1, 2, \ldots, k \}$. A quasi-state is added to $Next$ only if it is
    not in $Visited$, and every extracted quasi-state is added to $Visited$ and never
    removed. Therefore, no quasi-state is extracted twice. When $P$ is extracted,
    lines~\ref{alg:qdc:t}--\ref{alg:qdc:t_end} compute $t =
    \hat{\Delta}_\epsilon^i(\hat{\Delta}_f(P, \{ \sigma \})) = \tau_i(P, \sigma)$ for
    every $\sigma \in \Sigma$ and add the corresponding entries. If $t$ has not been
    extended before, lines~\ref{alg:qdc:e}--\ref{alg:qdc:e_next} compute $e = t \cup
    \hat{\Delta}_\epsilon^i(x) = \eta_i(t, x)$ for every $x \subseteq X$, add the
    corresponding entries, and add $e$ to $Next$ unless it has already been extracted.
    Otherwise, these entries have been added and these quasi-states have been added to
    $Next$ before. Therefore, the extracted quasi-states are exactly those obtained from
    $s_i = s^i$ by applications of $P \mapsto \eta_i(\tau_i(P, \sigma), x)$, that is,
    the elements of $U_i$; the extended quasi-states are exactly the elements of $T_i$;
    and the computed tables of component $i$ are those of $\mathcal{A}''$. So is the
    table of $\eta_x''$ computed at line~\ref{alg:qdc:etax}.
  \end{proof}

  \begin{proposition}\label{prop:quasidet_components_complexity}
    Let $\mathcal{A} = \langle \Sigma, Q, X, S, F, \Delta \rangle$ be a
    $k$-separated finite-state automaton. Then \Cref{alg:quasi_det_sfsa_components}
    completes on input $\mathcal{A}$ in $\mathcal{O}((|\Sigma| + 2^{|X|}) |Q| 2^{|X|}
    \sum_{i=1}^k 2^{|Q_i|})$ time.
  \end{proposition}

  \begin{proof}
    As in the proof of \Cref{prop:quasidet_complexity}, we assume that every operation
    on subsets of $Q$, including table look-ups and insertions, takes $\mathcal{O}(|Q|)$
    time. Line~\ref{alg:qdc:init} requires $\mathcal{O}(|Q|)$ time.
    Line~\ref{alg:qdc:etax} computes $4^{|X|}$ unions and requires
    $\mathcal{O}(2^{|X|} 2^{|X|} |Q|)$ time.
    Lines~\ref{alg:qdc:comp_begin}--\ref{alg:qdc:comp_init_end} are executed $k$ times
    and require $\mathcal{O}(|Q|)$ time.
    By \Cref{prop:quasidet_components_algorithm}, the body of the while loop is executed
    $|U_i| \leq 2^{|Q_i| + |X|}$ times for component $i$, and each execution requires
    $\mathcal{O}(|\Sigma| |Q|)$ time outside of
    lines~\ref{alg:qdc:e}--\ref{alg:qdc:e_next}. These lines are executed for $|T_i|
    \leq 2^{|Q_i| + |X|}$ quasi-states $t$ of component $i$, and require
    $\mathcal{O}(2^{|X|} |Q|)$ time for each of them. Therefore, the total runtime
    complexity bound for \Cref{alg:quasi_det_sfsa_components} is
    \begin{align*}
      &\mathcal{O}(2^{|X|} 2^{|X|} |Q| + k |Q| + (|\Sigma| + 2^{|X|}) |Q|
      \sum_{i=1}^k 2^{|Q_i| + |X|})\\
      &= \mathcal{O}((|\Sigma| + 2^{|X|}) |Q| 2^{|X|} \sum_{i=1}^k 2^{|Q_i|}),
    \end{align*}
    since $k \leq \sum_{i=1}^k 2^{|Q_i|}$ and $2^{|X|} \leq 2^{|X|} \sum_{i=1}^k
    2^{|Q_i|}$.
  \end{proof}

  \begin{theorem}\label{thm:quasidet_components}
    For every $k$-separated finite-state automaton $\mathcal{A} = \langle \Sigma, Q, X,
    S, F, \Delta \rangle$ a $k$-separated quasi-deterministic finite-state automaton
    $\mathcal{A}''$ with $\mathcal{L}_{\mathcal{A}''} = \mathcal{L}_\mathcal{A}$ and at
    most $2^{|Q_i| + |X|}$ quasi-states in its $i$-th component can be constructed in
    $\mathcal{O}((|\Sigma| + 2^{|X|}) |Q| 2^{|X|} \sum_{i=1}^k 2^{|Q_i|})$ time.
  \end{theorem}

  \begin{proof}
    Follows from \Cref{lem:quasidet_components} and
    \Cref{prop:quasidet_components_algorithm,prop:quasidet_components_complexity}.
  \end{proof}

  Unlike \Cref{prop:quasidet_complexity}, this bound does not depend on $D$. For the
  finite-state automata $\mathcal{A}_k$ above, $|U_i \cup T_i| \leq 4$, and
  \Cref{alg:quasi_det_sfsa_components} completes in time polynomial in $k$, whereas
  \Cref{alg:quasi_det_sfsa} executes at least $2^k$ iterations. If the size of the
  separator is bounded by a constant, the bound becomes $\mathcal{O}(|\Sigma| |Q|
  \sum_{i=1}^k 2^{|Q_i|})$. This is the case for finite-state automata constructed from
  regular expressions.

  \begin{corollary}\label{cor:quasidet_regex}
    Let $\mathcal{E}$ be a regular expression over an alphabet $\Sigma$ and let
    $\mathcal{A} = \langle \Sigma, Q, S, F, \Delta \rangle$ be the finite-state
    automaton constructed from $\mathcal{E}$ by the algorithm described in
    \Cref{thm:thompson}. Then a $2$-separated quasi-deterministic finite-state automaton
    $\mathcal{A}''$ such that $\mathcal{L}_{\mathcal{A}''} = \mathcal{L}_\mathcal{E}$
    can be constructed in $\mathcal{O}(|\Sigma| |Q| 2^{\frac{2}{3} |Q|})$ time.
  \end{corollary}

  \begin{proof}
    By \Cref{thm:thompson_separator}, \Cref{alg:thompson_separator} finds a separator
    $X$ and a partition $\{ X, Q_1, Q_2 \}_Q$ with $|X| = 2$ and $|Q_1|, |Q_2| \leq
    \frac{2}{3} |Q|$ in $\mathcal{O}(|Q|)$ time. By \Cref{thm:quasidet_components}
    applied with $k = 2$, \Cref{alg:quasi_det_sfsa_components} constructs
    $\mathcal{A}''$ with $\mathcal{L}_{\mathcal{A}''} = \mathcal{L}_\mathcal{A} =
    \mathcal{L}_\mathcal{E}$ in
    \begin{equation*}
      \mathcal{O}((|\Sigma| + 4) |Q| \cdot 4 \cdot (2^{|Q_1|} + 2^{|Q_2|})) =
      \mathcal{O}(|\Sigma| |Q| 2^{\frac{2}{3} |Q|})
    \end{equation*}
    time.
  \end{proof}

  In comparison, the deterministic finite-state automaton constructed by
  \Cref{thm:rabin_scott} can have up to $2^{|Q|}$ states, and \Cref{alg:quasi_det_sfsa}
  executes one iteration per state (\Cref{prop:quasidet_iterations}).
